% s-081-p081-m002; original PDF 81
\Aphorism{s-081-p081-m002}{74}
Dass die Natur nach menschlichen Begriffen unvernünftig und unmoralisch ist, bleibt noch ihre anständigste Eigenschaft: die einzige, die den Pantheismus halbwegs rechtfertigt. Freilich darf man, um sie zu würdigen, nicht selber zu den Schlechtweggekommenen, zu den \Foreign{in malam partem} unvernünftig und unmoralisch Behandelten gehören, — wie ja auch die Bewunderung der sogenannten „sittlichen Weltordnung“ nicht von Solchen getheilt wird, denen es dabei an den Kragen geht. In beiden Fällen muss man uninteressirt sein, um ästhetisch schätzen zu können, sei es, dass man die vorurtheilsfreie, natürliche Natur oder die menschlich gebundene, viereckige, vernünftig und moralisch missverstandene Natur höher schätzt.
\par

